\documentclass[12pt]{article}

\usepackage[T1]{fontenc}
\usepackage[margin=1in]{geometry}
\usepackage{amsmath,amssymb}
\usepackage{natbib}
\usepackage{hyperref}

\title{Introduction Roughs for SPECTRUM}
\author{}
\date{Drafted 2026-09-13; finalized after three rounds of adversarial review}

\begin{document}
\maketitle

\noindent This file contains seven candidate Introductions for
\texttt{01\_slr\_forecast\_intervention2026.tex}. Version 0 was written last,
in a separate session, as a short introduction that carries the essential
material and leaves the supporting detail of Versions 1--6 to the body of the
paper. It was subsequently revised for the audience of the AGU Chapman
Conference on sea-level rise projection, which raised its citation density and
its technical specificity. A later pass cut most of the West Antarctic
discussion, on the grounds that this study represents that uncertainty through
a scenario mixture rather than resolving it, and reinvested the space in the
central argument, which is the value of a hierarchy of forecast models that can
each be calibrated and validated against data. A further pass recast what
remained of the process-model discussion in general terms. It replaced the
catalogue of Antarctic process detail with the field-wide evidence that these
models are underconstrained relative to the observational record, hindcast
poorly, carry ensemble spreads that are not reliable uncertainty estimates,
depend on initialization, and depend on calibration choices as much as on
physics. That pass also replaced the non sequitur that followed the AR6
projection range, which had pivoted without warning into ISMIP6, LARMIP-2, and
structured expert judgment, with a statement of how wide the published high-end
range now is and of how the underlying ice-sheet models perform against the
historical record. West Antarctica is now named only once, in the closing
paragraph, and only to explain why that contributor is represented as a
scenario mixture. The general evidence retained in the model-class paragraph is
drawn in part from Antarctic studies but is framed as evidence about
sea-level projection as a whole. That pass dropped
ten keys (\texttt{Bamber2019}, \texttt{Edwards2019}, \texttt{Edwards2021},
\texttt{Gudmundsson2013}, \texttt{HaseloffSergienko2018},
\texttt{Levermann2020}, \texttt{Martin2026}, \texttt{Pegler2018},
\texttt{Ranganathan2024}, and \texttt{Seroussi2020}) and added none.

\noindent A final pass supplied references for the hierarchy-of-models
paragraph, which had been an uncited block of argument and would not have been
publishable as one. It now carries \texttt{Held2005}, \texttt{Jeevanjee2017},
and \texttt{Maher2019} for the model-hierarchy practice itself,
\texttt{Oreskes1994} for the point that a model of a natural system cannot be
verified against truth, \texttt{Aschwanden2021} and \texttt{Felikson2025} for
the formal role of observations in refining projections, \texttt{Lipscomb2025}
for the weight carried by agreement among models that do not share assumptions,
and \texttt{Green2015} with \texttt{Makridakis2020} for the value of simplicity,
whose published definition includes being understandable by those who must act
on the forecast. Version 0 now runs to 85 body lines, or 1.85 pages as typeset
at 46 lines per page, and cites 48 keys. Of those, \texttt{Lipscomb2025} was
verified against a local Zotero PDF and added to \texttt{references.bib}, and
\texttt{Jeevanjee2017}, \texttt{Maher2019}, and \texttt{Oreskes1994} were added
as new entries reconstructed from recall rather than from a publisher record or
a local PDF. Those three are flagged in \texttt{references.bib} and every field
in them must be confirmed before submission, as must \texttt{Held2005}, which
was already present and carries the same flag. Every other key is either
already used in the manuscript itself or in Versions 1--6 of this file, or was
verified against a local Zotero PDF (\texttt{Goldberg2026} and
\texttt{Morlighem2024}).
Versions 1--6 were produced as directed:
Version 1 follows the six-item bulleted outline already present at the top
of that manuscript's \texttt{\textbackslash section\{Introduction\}}; Version
2 is an independently constructed introduction drawing on the broader
survey in \texttt{literature\_component\_slr.tex} and the project's own
findings; Version 3 is a condensed synthesis of the two, retaining what we
judged the most load-bearing points (the stakes, the three-way model-class
trade-off, the WAIS tail-risk gap, and the SPECTRUM framework itself)
while dropping supporting detail and secondary citations; Versions 4 and 5
are two independent critical edits of Version 3 (described further below)
that substantially expand its treatment of West Antarctic Ice Sheet (WAIS)
uncertainty and restructure its sentences for classic-style, low-embedding
prose; Version 6 merges Versions 3--5 into a single text and corrects a
logic error common to all three (described further below). All six are
written to match the voice, quantitative style, and citation density of
the target manuscript. Versions 1 and 2 were each
drafted, then put through three rounds of independent adversarial review
(fresh reviewers each round, checking citation validity against
\texttt{references.bib}, accuracy against the papers actually cited, style
conformance, and prose quality) with a full revision pass after each
round, and were confirmed in the final round to compile cleanly (verified
with a live \texttt{pdflatex}/\texttt{bibtex} run) with no outstanding
material issues. Version 3 draws only on claims and citations already
verified in Versions 1 and 2, and received one adversarial review pass of
its own (citation validity, accuracy, and length target) rather than the
full three rounds, followed by a second revision pass (described below)
and, in Versions 4 and 5, a third round that substantially expands the
WAIS discussion with newly sourced, independently verified citations.
Nothing here has been copied into the manuscript itself.

\noindent\textbf{Length note.} Version 1 is 1044 words and Version 2 is
872 words (body text only, excluding headings, counted programmatically).
Version 3 originally targeted the intersection of ``half to two-thirds''
of each, i.e.\ 522--581 words, and its first draft came in at 580 words,
inside that band relative to both (0.56$\times$ V1, 0.67$\times$ V2). A
second revision pass (described below) added connective sentences to fix
structural and clarity problems surfaced by adversarial review; the
current text is 623 words (0.60$\times$ V1, 0.71$\times$ V2), modestly
above the original band, a trade-off made deliberately in favor of
clarity over the length target. Its citation count (23) is unchanged by
that pass and was scaled down roughly in proportion to its length rather
than held to the 25--40 range specified for Versions 1 and 2, since
forcing that many citations into a passage this short would come at the
expense of prose quality. Version 3's first draft received one full
adversarial review pass, which caught (and that pass fixed) several
compression artifacts: a fabricated causal link between the
model-parsimony citations and the held-out-validation design, an
inaccurate collapsing of Rahmstorf (2007) and Vermeer \& Rahmstorf (2009)
into a single-coefficient description (the latter has two), a dropped
hierarchy-of-models rationale, an orphaned citation, and a reversed
description of what component-level models trade away.

\noindent\textbf{Revision note (second pass).} The revised draft was then
checked by three independent adversarial reviewers, each rereading the
full text and mapping every sentence to the idea it is meant to convey,
one focused on argumentative structure and logical coherence, one on
plain-language and sentence-level clarity, and one on whether each
sentence's content actually matched the goal that mapping assigned it.
All three independently flagged the same unresolved referent (the phrase
``that gap'' opening Paragraph 3, which pointed to nothing named in the
preceding paragraph); two flagged a component-level-models sentence that
conflated a stated weakness (assuming stationarity) with an unrelated,
more important claim (no published implementation withholds the
aggregate record for an independent check) behind a single semicolon; and
reviewers separately flagged an unearned closing claim in Paragraph 2 (that
ordering model classes by complexity lets simpler ones check more complex
ones, asserted without the paragraph showing how), an unsupported closing
claim in Paragraph 3 (an appeal to ``adaptation planning'' introduced with
no groundwork), and several sentences overloaded with two or more distinct
claims joined by colons or semicolons. The text above fixes all of these:
it names the referent explicitly (``this missing independent check''),
splits the conflated and overloaded sentences, ties the Paragraph 2
closing claim to the specific held-out GMSLR--GMST check performed in
Paragraph 4, and rewrites the Paragraph 3 closing claim to point back to
the flooding-damage stakes already established in Paragraph 1 rather than
to an unintroduced concept. No citations were added or removed in this
pass.

\noindent\textbf{Revision note (third pass: Versions 4 and 5).} Two
independent editors were each given the post-second-pass text of Version 3
and the same brief: cut any sentence not earning its place, substantially
expand the WAIS paragraph with additional citations supporting the claim
that WAIS is the largest single source of sea-level-projection uncertainty
(not limited to a pre-supplied citation pool; both were explicitly told to
search the web for further real, verifiable literature), and restructure
sentences per the classic-style and syntax-tree principles in Steven
Pinker's \emph{The Sense of Style} (subject and verb kept close together,
complexity pushed to the edges of a sentence rather than center-embedded,
end-weighting, avoiding nominalizations, one claim per main clause).
Working independently, one editor (Version 4) built a thesis-first,
explicitly numbered four-line evidentiary case (observational trend,
independent expert judgment, multi-model spread, and a temporal-variance
decomposition), closing with a separate mechanism paragraph; the other
(Version 5) built a three-movement prose argument (mechanism, observation,
convergence) that additionally surfaces a genuine counter-finding
(\citealt{Edwards2019} finds ice-cliff collapse is not required to match
the observational record, bounding the 95th-percentile Antarctic
contribution below 43~cm without it) and argues this makes West Antarctic
uncertainty structural rather than narrowing it. Every new citation both
editors introduced was checked against a primary source (the publisher
page, a DOI resolver, or the Crossref API) before use, not recalled from
training data; four candidate citations were considered and rejected as
redundant with citations already used. The coordinating session
independently re-verified the four most load-bearing new numeric claims
(the \citealt{DeContoP2016} ``more than a meter by 2100'' figure, the
\citealt{Edwards2019} 43~cm bound, the \citealt{Seroussi2020}
$-$7.8~to~+30.0~cm ISMIP6 range, and the \citealt{Milillo2019}
tenfold spatial variability in Thwaites melt rates) against their
original abstracts via web search and found all four accurate as stated
in both drafts.

\noindent\textbf{Revision note (fourth pass: Version 6).} Versions 3, 4,
and 5 all opened the WAIS discussion with some variant of ``the stakes of
[the missing independent check] are concentrated in West Antarctica.''
This was flagged as a logic error: it frames West Antarctica's dominance
as a consequence of, or contingent on, the field-wide calibration gap
identified two paragraphs earlier, when in fact West Antarctica would be
the largest source of sea-level-projection uncertainty regardless of
whether any model class withholds an independent validation target. The
two facts are real but independent, and stringing them together with
``the stakes of X are concentrated in Y'' asserts a causal or scoping
relationship between them that the surrounding text never establishes.
Version 6 merges the strongest elements of Versions 3--5 (Version 4's
thesis-first framing and tighter model-class paragraph; Version 5's
Edwards2019 counter-finding, Milillo2019 and Levermann2020 citations, and
punchier closing sentences; the full citation union of both WAIS
expansions, with none dropped) and repairs the transition by (i) stating
West Antarctica's dominance as a freestanding claim, explicitly denying
that it follows from the calibration gap (``That is not a consequence of
how any model class chooses to calibrate; it follows from the ice
physics and the observational record themselves''), and (ii) moving the
correct, oppositely-directed causal claim to the start of the SPECTRUM
paragraph, where it belongs: the calibration gap matters \emph{because}
West Antarctica's stakes are largest, not the reverse (``the gap matters
most exactly where the stakes are largest: West Antarctica''). No
citations were added beyond the union already present in Versions 4 and
5, and none were dropped.

\noindent\textbf{Bibliography note.} Version 1 cites only keys that
already existed in \texttt{references.bib} at the time it was drafted.
Version 2 additionally cited several papers discussed in
\texttt{literature\_component\_slr.tex}: \texttt{Felikson2025},
\texttt{Slangen2022}, \texttt{Garner2018}, \texttt{Nauels2017},
\texttt{Kopp2023}, \texttt{McCormack2026}, \texttt{Bakker2017},
\texttt{Darnell2025}, and \texttt{Held2005}; these have since been added to
\texttt{references.bib} (by an earlier session, prior to the WAIS-expansion
pass described above), so all citations in Versions 1--3 now resolve.
Versions 4 and 5 introduce four further keys, added to
\texttt{references.bib} as part of this pass: \texttt{Kopp2017},
\texttt{Bulthuis2019}, \texttt{Milillo2019}, and \texttt{Levermann2020}.
Like \texttt{McCormack2026} and \texttt{Held2005} before them, these four
are flagged in \texttt{references.bib} with a note that they were
confirmed against the publisher record or the Crossref API rather than a
local Zotero PDF, and should be checked against a local PDF before
submission, per this project's citation-verification standard. Both new
sections also reuse the existing \texttt{DeContoP2016} and
\texttt{Edwards2019} keys (already present in \texttt{references.bib} and
already cited in Version 1) rather than introducing duplicate entries for
the same two papers.

\noindent\textbf{Bibliography note (fifth pass: Nerem2022).} A fifth key,
\texttt{Nerem2022} (\citealt{Nerem2022}: a purely empirical quadratic
extrapolation of the satellite-altimetry GMSL record itself, with no
temperature regression, found to agree with IPCC SROCC and AR6 21st-century
projections), was added to \texttt{references.bib} from a local Zotero PDF
and cited in all six versions above, in each version's semi-empirical
paragraph, as a still simpler empirical benchmark alongside
\texttt{Rahmstorf2007}/\texttt{Vermeer2009}: unlike those, it needs no
GMST regression at all, extrapolating the observed sea-level record
directly. No existing sentences were otherwise altered.

%% ====================================================================
\section*{Version 0 --- short introduction written to the manuscript's voice}


Every centimeter of that rise exposes roughly 1.6 million more people to annual coastal
flooding \citep{Kulp2019}, and without further adaptation annual flood damages reach
US\$14--18 trillion per year at 0.9--1.1~m of rise \citep{Jevrejeva2018}. What can be
done about that depends on how skillfully we can predict the rate of sea-level rise. Forecasts published over four decades span a wide range of methods and outcomes, and the
high-end estimates published since the Sixth Assessment Report of the Intergovernmental
Panel on Climate Change (AR6) span 0.6--2.0~m of rise by 2100 \citep{Garner2026}. Coastal
planning leans on AR6 itself, which projects 0.4--0.7~m this century
\citep{Fox-Kemper2021,Hirschfeld2023}. The ice-sheet models underlying those land-ice
projections reproduce the historical record poorly. Few match observed mass loss, and the
spread across an ensemble of model outputs is not a reliable measure of the uncertainty it
is taken to represent \citep{Aschwanden2021}. 

Global warming that has happened and will happen ensure that projections of sea-level rise for the 21st century and beyond are out of sample with respect to the modern observational record no matter the projection method. The central difficulty is therefore not building a model but calibrating it, testing it against data
withheld from that calibration, and reporting uncertainties in a form that separates the
contribution of each physical process. A recent community synthesis makes this
case, calling for projections narrowed by observations and for uncertainty characterized
by its structure and not only by its magnitude \citep{Felikson2025}.

Sea-level projection models span a wide range of complexity, and their strengths are
complementary. Process-based ice-sheet models solve the equations of ice flow,
deformation, and fracture, and they are where competing hypotheses about the mechanisms
of ice loss can be posed against one another and tested \citep{Schoof2007,Morlighem2024}.
Their parameters, such as the basal sliding coefficient, are inferred from observations by ill-posed inverse problems in which the data constrain only a limited set of parameter directions \citep{Isaac2015,Petra2014}. They depart from reconstructions of historical GMSLR in hindcast \citep{Aschwanden2021}, and most
projections initialized in 2016 grow the Antarctic Ice Sheet through the years that
overlap the satellite record, with outcomes that remain sensitive to the state a model starts from \citep{Fricker2025}. Structural assumptions carry comparable weight. The stress exponent of the flow law for ice is usually assigned a value lower than observations of fast-flowing ice shelves indicate \citep{Millstein2022}, and where grounding lines retreat rapidly, uncertainty in that exponent produces a wider spread in projected ice loss than uncertainty in climate forcing does \citep{GetraerMorlighem2025}. Calibration can hide such a bias, because a model that assumes $n = 3$ can be initialized to match observations of an $n = 4$ glacier yet underestimate its 100-year sea-level contribution by 21 to 35\% in retreat benchmarks \citep{Martin2026}. Two ice-sheet models calibrated against different observations of the same glacier give mid-century loss rates differing by more than a factor of two \citep{Goldberg2026}. %Forecasting research across many fields finds that models carrying more free parameters than their data can constrain forecast worse than simpler ones \citep{Green2015,Makridakis2020}.

At the other extreme, extrapolations of the observed record and semi-empirical fits of
the rate of GMSLR to global mean surface temperature (GMST) are transparent and anchored
in the quantity being predicted \citep{Nerem2022,Rahmstorf2007,Vermeer2009,Kopp2016}, and
they track the observed record at least as well as contemporaneous process-based
projections \citep{Rahmstorf2012}. They give up any account of which processes are
responsible, and the fitted relationship need not hold as individual contributors respond
on different timescales \citep{Church2013}. Between the two sit component-level models,
which separate GMSLR into its physical contributors and calibrate each one against
observations, model output, or expert assessment
\citep{Kopp2014,Mengel2016,Wong2017,Perrette2013,Nauels2017,Kopp2023}, trading spatial
detail for parsimony and a direct link to data.

A numerical model of a natural system can be tested against data and against other
models, but it cannot be shown to be correct \citep{Oreskes1994}. No one class is at once
the most physically detailed, the most transparent, and the most tightly constrained by
data, so treating any one of them as authoritative discards what the others offer.
Ordered by complexity they instead form a scaffolding, a practice long established
elsewhere in climate science, where hierarchies spanning idealized and comprehensive
models have served understanding better than comprehensive models alone
\citep{Held2005,Jeevanjee2017,Maher2019}. Each rung can be calibrated and tested against data before the next is
trusted, giving observations the formal role in refining projections that recent
assessments of ice-sheet and sea-level projection have called for
\citep{Aschwanden2021,Felikson2025}. An extrapolation of
the observed record sets a scenario-independent reference that any more elaborate model
should be able to reproduce. A semi-empirical fit adds a dependence on warming and little
else, which isolates that sensitivity and makes it testable on its own. A
component-level model then asks whether the same sensitivity can be recovered from the
separate physics of the ocean, the glaciers, and the ice sheets, each calibrated against
its own record. A process-based model asks which mechanisms within those contributors are
responsible. Calibration and validation are then carried out within the space of models
rather than against a single preferred one, and agreement between rungs becomes a
diagnostic in its own right. Where the rungs agree, a result rests on more than one set
of assumptions, and agreement of that kind carries more weight than agreement among
models that share their assumptions \citep{Lipscomb2025}. Where they diverge, the
divergence localizes what the data do not yet
constrain and identifies the measurement that would settle it. The simplest rungs often
carry the most transferable insight, because what they assume is small enough to state
plainly, to test directly, and to be understood by those who must act on it
\citep{Green2015,Makridakis2020}. A scaffolding built this way complements process-based
modeling rather than substituting for it.

Here we construct a data-driven, component-level forecast of 21st-century sea-level
rise. Where the observational record can constrain a contributor to GMSLR, we give it a
parsimonious physical model calibrated by Bayesian inference against an independent
record spanning decades, among them NOAA thermosteric estimates
\citep{Levitus2012}, EN4 subsurface ocean temperature \citep{Good2013}, Greenland
outlet-glacier discharge \citep{Mouginot2019}, GlaMBIE \citep{GlaMBIE2025}, and IMBIE-3
\citep{Otosaka2026}. The remaining contributors, among them Greenland and East Antarctic
surface mass balance and terrestrial water storage, are adopted from published model
sensitivities and assessments \citep{Fettweis2020,Noel2021,Frieler2015,Frederikse2020}.
We call the framework SPECTRUM, for Sea-level Projections from Empirical Calibration of
Thermodynamical, Reduced-physics Models. Most of its component models are adapted from
formulations in the literature, and for those our contribution lies in the calibration
data, the validation, and the treatment of uncertainty. The exception is Greenland Ice
Sheet discharge, for which we develop a three-parameter model driven by the temperature
of the surrounding ocean. West Antarctic discharge is the one contributor we do not fit
to warming, because the observational record cannot yet distinguish a self-sustaining
retreat from one sustained by ocean melting from below
\citep{Sergienko2026,Hill2023,Reese2023}, and we represent it instead as a structured
mixture of scenarios rather than a single trajectory. We do not attempt to settle that
question here. Observed GMSLR and its sensitivity to observed GMST are withheld from
every component calibration and reserved as diagnostics for the summed model, so the
central test of the framework is out of sample. The independently calibrated components
close the satellite altimetry budget \citep{Beckley2017,Hamlington2020} and reproduce the
level, the rate, and the transient sensitivity of GMSLR since 1950 across roughly
1 $^\circ$C of warming \citep{Frederikse2020,Fox-Kemper2021}, a record none of them ever
saw.

%% ====================================================================
\section*{Version 1 --- follows the manuscript's own itemized outline}

Coastal communities plan against a moving target. Each additional
centimeter of global mean sea-level rise (GMSLR) exposes roughly 1.6
million more people to annual flooding \citep{Kulp2019}, nearly 90\% of
them in low- and middle-income countries \citep{Jongman2012}. Without
further adaptation, annual flood damages are projected to reach
US\$14--18T (trillion) per year once GMSLR reaches 0.9--1.1~m
\citep{Jevrejeva2018}, while a lower-bound estimate of the cost of adapting
to that rise, based on the price of raising dikes along a small fraction of
the world's coastlines, is on the order of US\$4T per meter
\citep{Tiggeloven2020}. Developing countries currently receive roughly
US\$21B per year in adaptation finance across all sectors
\citep{UNEP2023}, far short of what these cost estimates imply is needed.
That shortfall makes the reliability of the sea-level projections used to
plan adaptation directly consequential for how effectively it can be
closed.

Sea-level projections have been published regularly for more than four
decades and span a wide range of methods and outcomes \citep{Garner2026}.
Rates of GMSLR have already doubled twice since the mid-20th century,
first over about 16 years and $0.34\ ^\circ$C of warming and again over
about 21 years and $0.41\ ^\circ$C \citep{Frederikse2020,Hamlington2024,
Rohde2020}. Yet current coastal planning still leans heavily on
projections from the IPCC Sixth Assessment Report (AR6), which places
21st-century rise at 0.4--0.7~m \citep{Fox-Kemper2021,Hirschfeld2023}.
This variety is not a weakness so much as an underused resource: no single
class of model is simultaneously the most physically detailed, the most
transparent, and the most tightly constrained by data, and a framework
that arrays these approaches by complexity, testing each rung against data
before trusting the next, offers a way to combine their respective
strengths and to identify where projections agree, diverge, or exceed what
the data can support \citep{Green2015,Makridakis2020}.

At the most detailed end sit numerical ice-sheet models, which resolve ice
flow, deformation, and fracture directly and let competing physical
hypotheses, such as whether marine ice-cliff instability can activate this
century, be posed and debated within a common numerical framework
\citep{DeContoP2016,Edwards2019,Morlighem2024}. Their principal weakness
is that this physical detail demands more constraints than the
observational record can currently supply. In hindcasting studies,
process-based models diverge from historical GMSLR reconstructions
\citep{Aschwanden2021} and from satellite observations within a decade of
initialization \citep{Fricker2025}. These disagreements trace to
structural gaps such as the omission of leading-order calving physics
\citep{Fricker2025} and an assumed ice viscosity far less sensitive to
stress than observations indicate \citep{Millstein2022,
GetraerMorlighem2025}, compounded by initialization uncertainty and, for
marine ice sheets specifically, dynamics that can skew the resulting
projections toward high outcomes rather than merely widening them
\citep{Robel2019}. The resulting parameterizations are numerous enough,
and their interactions complex enough, that these models function largely
as black boxes even to their developers, and every century-scale
projection necessarily extrapolates them well beyond the short
observational record used to tune them. Decades of general forecasting
research show that
models with more free parameters than their calibration data can
constrain tend to forecast worse, not better, than simpler alternatives
\citep{Green2015}.

At the opposite extreme are semi-empirical models that regress the
observed rate of GMSLR directly onto global mean surface temperature
(GMST), a class whose transparency and direct grounding in the data being
predicted make it easy to test and to communicate \citep{Rahmstorf2007,
Vermeer2009}. Statistically, semi-empirical fits track the observed
sea-level record over multi-decade hindcast windows at least as well as,
and at times better than, contemporaneous process-based projections
\citep{Rahmstorf2012}, and
Bayesian treatments of the same GMST-GMSL relationship, which propagate
measurement and proxy uncertainty explicitly rather than fitting a single
point estimate, extend the approach over far longer records
\citep{Kopp2016}. Their central limitation is that a single global
regression coefficient is not obviously physical or causal, may not stay
stationary as individual components respond to different forcings on
different timescales, and gives no way to tell, when it forecasts
correctly, whether that agreement reflects the right physical reasons
\citep{Church2013}. An even simpler benchmark needs no temperature
regression at all: extrapolating a quadratic fit to the
satellite-observed GMSL record itself agrees with IPCC SROCC and AR6
projections to within their own stated uncertainty \citep{Nerem2022}.

Between these extremes sit component-level models that decompose GMSLR
into its physical contributors, ocean thermal expansion, glaciers, and the
ice sheets, and calibrate each against existing observational or model
data \citep{Kopp2014,Mengel2016,Wong2017}. This intermediate complexity is
a deliberate trade. Parsimony and a direct connection to physical
processes come at the cost of lumped parameters that cannot resolve
individual basins or glaciers, an assumption of stationarity over the
calibration period, and the same extrapolation risk that limits every
other class of model. The calibration strategy within this class still
varies: some scale a component decomposition to modeled temperature
trajectories \citep{Perrette2013}, others use Bayesian-calibrated
reduced-complexity models, with an expert prior on the most uncertain term
(Antarctic dynamics), to bound a physically plausible upper limit on total
sea-level rise \citep{Sriver2012}, and still others forgo new
component physics altogether and instead statistically fuse existing
probabilistic projections into a single consensus estimate
\citep{Grandey2024}.

A framework that spans this hierarchy, from semi-empirical regression
through component-level, data-calibrated models to full process-based
simulation, does not need to choose a single class as authoritative. It can use agreement
or disagreement across rungs as a diagnostic in its own right, letting the
simplest, most transparent models motivate and check the more complex
ones, calibrating and validating within the space of models rather than
against a single preferred one, and it can let the most consequential
insights come from the simplest rung rather than the most complex. Here,
we construct such a framework:
a data-driven, component-level sea-level projection that treats each
contributor to GMSLR, ocean thermal expansion, Greenland Ice Sheet (GrIS)
surface mass balance (SMB) and glacier discharge, West Antarctic Ice Sheet
(WAIS) glacier discharge, East Antarctic Ice Sheet (EAIS) SMB, the
Antarctic Peninsula, glaciers outside the ice sheets, and terrestrial water
storage, with its own parsimonious physical model, calibrated against
independent observational records using Bayesian inference. We call this
framework SPECTRUM, for Sea-level Projections from Empirical Calibration
of Thermodynamical, Reduced-physics Models. Most of its component models
adapt existing formulations from the literature; the exception is
Greenland Ice Sheet discharge, whose 3-parameter model connecting glacier
discharge directly to ocean temperature is developed here. Critically,
observed GMSLR and its sensitivity to observed GMST are withheld from
every component-level calibration and used only as diagnostics on the
composite model, so that the framework's central test, whether
independently calibrated parts sum to reproduce a whole they were never
fit to, is genuinely out of sample.

%% ====================================================================
\section*{Version 2 --- independent synthesis from the literature review}

Every additional centimeter of global mean sea-level rise (GMSLR) exposes
roughly 1.6 million more people to annual coastal flooding
\citep{Kulp2019}, and unmitigated damages are projected to reach
US\$14--18T (trillion) per year once rise reaches 0.9--1.1~m
\citep{Jevrejeva2018}. Sea-level projections meant to guide adaptation
against those costs have been published for more than four decades. A
database of 756 such projections drawn from 103 studies shows their
production concentrated in a handful of wealthy countries and their
horizons only recently extended past 2100, and finds that, prior to the
IPCC's Sixth Assessment Report (AR6), individual studies' high-end
estimates routinely exceeded the IPCC's own central ranges
\citep{Garner2026}. Over that same period, successive IPCC assessment
cycles have revised their component treatments and uncertainty structure
substantially \citep{Slangen2022,Garner2018}, evidence less of an immature
field than of one still searching for the right relationship between
physical detail and the observational record available to constrain it.

A growing, cross-institutional view holds that the way forward is not more
physical detail but tighter observational constraint. A recent synthesis
spanning the leading process-based and observational sea-level research
groups argues explicitly for narrowing projections with data and for
characterizing uncertainty structure, not just magnitude, in both the
observations and the models \citep{Felikson2025}. That argument responds
to a specific, well-documented weakness of the process-based ice-sheet
models that dominate current assessments, including AR6's own 0.4--0.7~m
projection for 21st-century rise \citep{Fox-Kemper2021}. In hindcasting
exercises, these models diverge from the historical GMSLR record
\citep{Aschwanden2021} and from satellite observations within a decade of
initialization \citep{Fricker2025}, disagreements traced to structural
gaps such as an assumed ice viscosity far less sensitive to stress than
observations indicate \citep{Millstein2022,GetraerMorlighem2025}.

Component-summation frameworks, which decompose GMSLR into its physical
contributors and combine them probabilistically, already offer a partial
answer, and existing implementations differ mainly in what each component
is calibrated against: expert-assessed probability distributions
\citep{Kopp2014}, independently calibrated semi-empirical response
functions \citep{Mengel2016}, jointly calibrated physical models
\citep{Wong2017}, or reduced-complexity emulators of process-model output
\citep{Nauels2017}, with the modular platform now underlying AR6 designed
to accommodate all of these calibration strategies within a common
structure \citep{Kopp2023}. What none of them do is hold out the aggregate
observational record, global sea level's own sensitivity to global
temperature, entirely from every component-level fit and use it solely as
an independent check on the summed result.

The stakes of that calibration choice are concentrated in West Antarctica.
Marine ice sheet instability does not simply widen the uncertainty on
Antarctic contributions; it mathematically skews the resulting
distribution toward high outcomes \citep{Robel2019}, and independent
regional ocean modeling finds that ocean-driven melt of West Antarctic ice
shelves accelerates at similar rates across emissions scenarios this
century, meaning much of that risk may already be locked in
\citep{Naughten2023}. Recent work also finds that the near-term rate of
Antarctic mass loss is itself predictive of the rate over the following
3--5 decades \citep{McCormack2026}, suggesting the tail is not
only large but partially foreseeable from the current observational
record. Yet deep-uncertainty treatments that represent this risk as a
structured scenario mixture rather than a single process-model trajectory
remain rare \citep[e.g.,][]{Bakker2017}. A recent variance decomposition
reaches a seemingly contrary conclusion, that emissions-scenario choice
dominates projection uncertainty through the mid-century before
geophysical uncertainty takes over \citep{Darnell2025}, a claim that a
framework treating the tail explicitly cannot simply set aside.

None of this is an argument against physical detail; it is an argument
for testing that detail against data before trusting it. Outside
sea-level rise specifically, climate science has long recognized that a
hierarchy spanning idealized and comprehensive models serves
understanding, and often forecast robustness, better than comprehensive
models alone \citep{Held2005}, a pattern general forecasting research
confirms empirically across some 100,000 time series
\citep{Green2015,Makridakis2020}. Semi-empirical models that regress GMSLR
directly onto GMST provide another genuinely data-driven benchmark
\citep{Rahmstorf2007,Vermeer2009}, distinguishing the choice of model
complexity from the choice of calibration source. An even simpler
benchmark needs no temperature regression at all: a quadratic
extrapolation of the satellite-observed GMSL record alone agrees with
IPCC SROCC and AR6 projections to within their own stated uncertainty
\citep{Nerem2022}. Nor is the record
uniformly unkind to assessment-based projection. An earlier IPCC
assessment's mid-range 21st-century sea-level projection verified to
within about a centimeter three decades later \citep{Tornqvist2025}, a
reminder that any critique of process-based projections should be aimed at
specific assessment cycles and specific components, not at process-based
reasoning as such.

Here, we build a component-level framework that puts that discipline into
practice directly: every physical contributor to GMSLR, ocean thermal
expansion, Greenland Ice Sheet (GrIS) surface mass balance and glacier
discharge, West Antarctic Ice Sheet (WAIS) glacier discharge, East
Antarctic Ice Sheet (EAIS) surface mass balance, the Antarctic Peninsula,
glaciers outside the ice sheets, and terrestrial water storage, is
calibrated with its own parsimonious physical model against an
independent observational record, using Bayesian inference to propagate
parameter uncertainty rather than adopting it from assessment or
process-model output. We call the result SPECTRUM, for Sea-level
Projections from Empirical Calibration of Thermodynamical, Reduced-physics
Models. The only genuinely new component model links GrIS discharge to
ocean temperature through a 3-parameter time-lagged response; every other
component adapts an existing formulation. Its treatment of WAIS is a
structured
deep-uncertainty scenario mixture built to represent the skewed tail
described above rather than average over it. Observed GMSLR and its GMST
sensitivity, withheld from every component calibration, are reserved as
the composite-level diagnostic that current frameworks do not report, and
against which SPECTRUM's aggregated components are tested here for the
first time.

%% ====================================================================
\section*{Version 3 --- condensed synthesis}

Global mean sea-level rise (GMSLR) already exposes roughly 1.6 million
more people to annual flooding for every additional centimeter
\citep{Kulp2019}, with unmitigated damages projected at US\$14--18T
(trillion) per year once rise reaches 0.9--1.1~m \citep{Jevrejeva2018}.
Sea-level projections meant to inform that risk have been published for
four decades and span a wide range of methods and outcomes
\citep{Garner2026}. Rates of GMSLR have already doubled twice since the
mid-20th century \citep{Frederikse2020,Hamlington2024}, yet coastal
planning still leans heavily on the IPCC Sixth Assessment Report (AR6),
which projects 0.4--0.7~m of 21st-century rise
\citep{Fox-Kemper2021,Hirschfeld2023}. A growing, cross-institutional view
holds that narrowing AR6's range requires tighter observational
constraint on projections, not more physical detail \citep{Felikson2025}.

Three model classes have taken up that challenge, each with a real
weakness. Process-based ice-sheet models resolve ice flow and
fracture directly, but they demand more constraints than the
observational record can currently supply. In hindcasting studies, they
diverge from the historical GMSLR record \citep{Aschwanden2021} and from
satellite observations within a decade of initialization
\citep{Fricker2025}. These disagreements trace to gaps such as an
assumed ice viscosity far less sensitive to stress than observations
indicate \citep{Millstein2022,GetraerMorlighem2025}. Decades of general
forecasting research show that models with more free parameters than
their data can constrain tend to forecast worse than simpler alternatives
\citep{Green2015,Makridakis2020}. Semi-empirical models that regress the
rate of GMSLR onto global mean surface temperature are simple and
transparent, but the underlying relationship is not obviously physical,
causal, or stationary \citep{Rahmstorf2007,Vermeer2009,Church2013}. Even
so, such fits track the observed record at least as well as
contemporaneous process-based projections \citep{Rahmstorf2012}. An even
simpler benchmark needs no temperature regression at all: a quadratic
extrapolation of the satellite-observed GMSL record alone agrees with
IPCC SROCC and AR6 projections to within their own stated uncertainty
\citep{Nerem2022}. Between
these extremes, component-level models decompose GMSLR into its physical
contributors and calibrate each against existing data
\citep{Kopp2014,Mengel2016,Wong2017}, trading basin-scale resolution for
parsimony and a direct physical connection, at the cost of assuming
stationarity. None of the published implementations, however, withholds
the aggregate observational record from every component fit, so none can
use it as an independent check on the summed result. Ordering these
three model classes by complexity, rather than treating any one as
authoritative, turns their disagreement into a diagnostic, letting the
simplest, most transparent models motivate and check the more complex
ones.

The stakes of this missing independent check are concentrated in West
Antarctica, where marine ice sheet instability does not merely widen
projection uncertainty but skews it toward high outcomes
\citep{Robel2019}, and where independent regional ocean modeling finds
melt accelerating at similar rates across emissions scenarios this
century, meaning much of that risk may already be locked in
\citep{Naughten2023}. This West Antarctic risk differs structurally from
the temperature-sensitive uncertainty in the other components. Treating
it as a single central trajectory, rather than a mixture spanning its
full range, would understate exactly the high-end damages introduced
above.

Here, we build a component-level framework in which each contributor to
GMSLR (ocean thermal expansion; Greenland Ice Sheet (GrIS) surface mass
balance and glacier discharge; West Antarctic Ice Sheet (WAIS) glacier
discharge; East Antarctic Ice Sheet (EAIS) surface mass balance; the
Antarctic Peninsula; glaciers outside the ice sheets; and terrestrial
water storage) gets its own parsimonious physical model, calibrated
against an independent observational record using Bayesian inference. We
call it SPECTRUM, for Sea-level Projections from Empirical Calibration of
Thermodynamical, Reduced-physics Models. Its only new component
model links GrIS discharge to ocean temperature through a 3-parameter
time-lagged response. Its treatment of WAIS is a structured
deep-uncertainty scenario mixture, built to represent the skewed tail
described above rather than average over it. Observed GMSLR and its
sensitivity to GMST, the relationship at the core of the semi-empirical
class described above, are withheld from every component calibration and
used only as a diagnostic on the composite model. The central test is
therefore out of sample. It asks whether the independently calibrated
parts reproduce a whole they were never fit to.

%% ====================================================================
\section*{Version 4 --- critically edited from Version 3 (draft A): thesis-first, four-line WAIS case}

Global mean sea-level rise (GMSLR) already exposes roughly 1.6
million more people to annual flooding for every additional centimeter
\citep{Kulp2019}, with unmitigated damages projected at US\$14--18T
(trillion) per year once rise reaches 0.9--1.1~m \citep{Jevrejeva2018}.
Sea-level projections meant to inform that risk have been published for
four decades and span a wide range of methods and outcomes
\citep{Garner2026}. Rates of GMSLR have already doubled twice since the
mid-20th century \citep{Frederikse2020,Hamlington2024}. Coastal planning
nonetheless leans heavily on the IPCC Sixth Assessment Report (AR6),
which projects 0.4--0.7~m of 21st-century rise
\citep{Fox-Kemper2021,Hirschfeld2023}. A growing, cross-institutional
view holds that narrowing AR6's range requires tighter observational
constraint on projections, not more physical detail \citep{Felikson2025}.

Three model classes have taken up that challenge, each with a real
weakness. Process-based ice-sheet models resolve ice flow and
fracture directly, but they demand more constraints than the
observational record can currently supply. In hindcasting studies, they
diverge from the historical GMSLR record \citep{Aschwanden2021} and from
satellite observations within a decade of initialization
\citep{Fricker2025}. These disagreements trace to gaps such as an
assumed ice viscosity far less sensitive to stress than observations
indicate \citep{Millstein2022,GetraerMorlighem2025}. Decades of general
forecasting research show that models tend to forecast worse than
simpler alternatives once they carry more free parameters than their
data can constrain \citep{Green2015,Makridakis2020}. Semi-empirical
models regress the rate of GMSLR onto global mean surface temperature.
They are simple and transparent, but the underlying relationship is not
obviously physical, causal, or stationary
\citep{Rahmstorf2007,Vermeer2009,Church2013}. Even so, such fits track
the observed record at least as well as contemporaneous process-based
projections \citep{Rahmstorf2012}. An even simpler benchmark needs no
temperature regression at all: a quadratic extrapolation of the
satellite-observed GMSL record alone agrees with IPCC SROCC and AR6
projections to within their own stated uncertainty \citep{Nerem2022}.
Between these extremes,
component-level models decompose GMSLR into its physical contributors
and calibrate each against existing data
\citep{Kopp2014,Mengel2016,Wong2017}, trading basin-scale resolution for
parsimony and a direct physical connection, while assuming the
relationships they calibrate are stationary. None of the published
implementations, however, withholds the aggregate observational record
from every component fit, so none can use it as an independent check on
the summed result. Rather than treat any one model class as
authoritative, we order the three by complexity and treat their
disagreement as a diagnostic. The simplest, most transparent models
motivate and check the more complex ones.

The stakes of this missing independent check are concentrated in West
Antarctica. Four independent lines of evidence converge on the same
conclusion. West Antarctica is the largest source of uncertainty in
sea-level projections, and its share of that uncertainty is growing.

First, the observational record shows the loss is already under way.
Antarctica's grounding lines have retreated for thirty consecutive years
\citep{Rignot2026}, and the Thwaites Glacier basin in West Antarctica
may already be undergoing an unstable, self-sustaining collapse
\citep{Joughin2014}, with recent observations there consistent with
continued high rates of loss over the next fifty years
\citep{Goldberg2026}. Second, a structured survey of expert judgment,
built independently of any process-model ensemble, assigns West
Antarctica wider and more heavily tailed contribution distributions than
those ensembles do \citep{Bamber2019}. Third, that spread is not merely
disagreement between models talking past each other. A single ice-sheet
model's own uncertainty in calving, basal sliding, ice-shelf rheology,
and ocean forcing, propagated forward for centuries, already produces a
wide range of Antarctic outcomes \citep{Bulthuis2019}, and the largest
multi-model ensemble, ISMIP6, adds still more disagreement among its
Antarctic members than among the members representing any other
component \citep{Seroussi2020}. Fourth, that disagreement is not fixed
in time. A variance decomposition of projected sea-level rise finds
emissions-scenario choice dominant through mid-century, after which
geophysical uncertainty, chiefly Antarctic ice dynamics, becomes the
largest term \citep{Darnell2025}, and most of the emissions-pathway
spread in total projections traces to the Antarctic contribution alone
\citep{DeConto2021}. One recent finding cuts the other way: the
near-term rate of Antarctic mass loss appears to carry predictive
information about the rate over the following three to five decades
\citep{McCormack2026}.

The physical mechanism behind West Antarctic loss argues against
reading that predictive signal as reassurance. Marine ice sheet
instability does not merely widen projection uncertainty. It skews the
distribution toward high outcomes \citep{Robel2019}. Coupled with
hydrofracturing and ice-cliff collapse, that instability could add more
than a meter of Antarctic sea-level rise by 2100 under unabated
emissions \citep{DeContoP2016}. Independent regional ocean modeling
finds shelf melt accelerating at similar rates across emissions
scenarios this century, meaning much of that risk is already locked in
regardless of future mitigation \citep{Naughten2023}. West Antarctic
risk therefore differs structurally from the temperature-sensitive
uncertainty in the other components. It is skewed, partly irreducible,
and growing in relative importance as the century proceeds, and even
the choice of which physical processes to include, not just which
parameters to assume, changes the shape of its probability distribution
\citep{Kopp2017}. Some published projections already treat it that way,
as a deep-uncertainty mixture of scenarios rather than a single central
trajectory \citep{Bakker2017}. Anything less would understate exactly
the high-end damages introduced above.

Here, we build a component-level framework that separates GMSLR into
seven contributors: ocean thermal expansion; Greenland Ice Sheet (GrIS)
surface mass balance and glacier discharge; West Antarctic Ice Sheet
(WAIS) glacier discharge; East Antarctic Ice Sheet (EAIS) surface mass
balance; the Antarctic Peninsula; glaciers outside the ice sheets; and
terrestrial water storage. Each contributor gets its own parsimonious
physical model, calibrated against an independent observational record
using Bayesian inference. We call it SPECTRUM, for Sea-level
Projections from Empirical Calibration of Thermodynamical,
Reduced-physics Models. Its only new component model links GrIS
discharge to ocean temperature through a 3-parameter time-lagged
response. Its treatment of WAIS is a structured deep-uncertainty
scenario mixture, built to represent the skewed tail described above
rather than average over it. We withhold observed GMSLR and its
sensitivity to GMST, the relationship at the core of the semi-empirical
class described above, from every component calibration, and use it
only as a diagnostic on the composite model. The central test is
therefore out of sample. It asks whether the independently calibrated
parts reproduce a whole they were never fit to.

%% ====================================================================
\section*{Version 5 --- critically edited from Version 3 (draft B): three-movement WAIS argument with counter-evidence}

Global mean sea-level rise (GMSLR) already exposes roughly 1.6
million more people to annual flooding for every additional centimeter
\citep{Kulp2019}, with unmitigated damages projected at US\$14--18
trillion per year once rise reaches 0.9--1.1~m \citep{Jevrejeva2018}.
Projections of that risk span four decades and a wide range of methods
and outcomes \citep{Garner2026}. Rates of GMSLR have already doubled
twice since the mid-20th century \citep{Frederikse2020,Hamlington2024},
yet coastal planning still leans heavily on the IPCC Sixth Assessment
Report (AR6), which projects 0.4--0.7~m of 21st-century rise
\citep{Fox-Kemper2021,Hirschfeld2023}. A growing, cross-institutional
view holds that narrowing AR6's range requires tighter observational
constraint on projections, not more physical detail \citep{Felikson2025}.

Three model classes have taken up that challenge, each with a real
weakness. Process-based ice-sheet models resolve ice flow and fracture
directly, but they demand more constraints than the observational
record can currently supply. In hindcasting studies, they diverge from
the historical GMSLR record \citep{Aschwanden2021} and from satellite
observations within a decade of initialization \citep{Fricker2025}.
These disagreements trace to gaps such as an assumed ice viscosity far
less sensitive to stress than observations indicate
\citep{Millstein2022,GetraerMorlighem2025}. Forecasting research
spanning decades and many fields finds the same pattern. Models with
more free parameters than their data can constrain tend to forecast
worse than simpler alternatives \citep{Green2015,Makridakis2020}.
Semi-empirical models that regress the rate of GMSLR onto global mean
surface temperature are simple and transparent, but the underlying
relationship is not obviously physical, causal, or stationary
\citep{Rahmstorf2007,Vermeer2009,Church2013}. Even so, such fits track
the observed record at least as well as contemporaneous process-based
projections \citep{Rahmstorf2012}. An even simpler benchmark needs no
temperature regression at all: a quadratic extrapolation of the
satellite-observed GMSL record alone agrees with IPCC SROCC and AR6
projections to within their own stated uncertainty \citep{Nerem2022}.
Between these extremes sit
component-level models, which decompose GMSLR into its physical
contributors and calibrate each against existing data
\citep{Kopp2014,Mengel2016,Wong2017}. This approach trades basin-scale
resolution for parsimony and a direct physical connection, at the cost
of assuming stationarity. None of the published implementations,
however, withholds the aggregate observational record from every
component fit, so none can check the summed result against it
independently. No one class is authoritative on its own. Ranked instead
by complexity, the three become a diagnostic. The simplest, most
transparent models motivate and check the more complex ones.

The stakes of that missing check are concentrated in West Antarctica.
Marine ice sheet instability there does not merely widen the
uncertainty in sea-level projections, it skews that uncertainty toward
high outcomes \citep{Robel2019}. A second, compounding process operates
there too. Hydrofracturing strips away the ice shelves that buttress
tall ice cliffs, letting them collapse under their own weight. The
model that first identified this mechanism found it could add more
than a meter to Antarctica's contribution by 2100 \citep{DeContoP2016}.
How much of that process will actually occur remains itself disputed.
Recalibrating the same physics against observations from the last
interglacial, the Pliocene, and the past 25 years found no need to
invoke ice-cliff collapse at all, and put the resulting 95th-percentile
Antarctic contribution below 43~cm this century \citep{Edwards2019}.
West Antarctica therefore carries a layer of uncertainty the other
components do not. Two things are unresolved there, not one: whether an
instability occurs at all, and how large it would be if it did.

Thwaites Glacier, the drainage basin most likely to host either
instability, may already be undergoing a self-sustaining retreat that
cannot be stopped, according to observations and modeling published as
early as 2014 \citep{Joughin2014}. Later radar surveys found that its
retreat proceeds unevenly along the grounding line, with warm water
undercutting some sectors ten times faster than others
\citep{Milillo2019}. More recent observations remain consistent with
continued high rates of loss there over the next 50 years
\citep{Goldberg2026}, and grounding lines across Antarctica have
retreated persistently for three decades \citep{Rignot2026}. Independent
regional ocean modeling adds a further finding. Ice-shelf melt
accelerates at similar rates whether emissions are cut sharply or not
this century, so much of this risk is already locked in regardless of
the mitigation path taken \citep{Naughten2023}.

Several other, methodologically independent assessments converge on the
same conclusion: West Antarctica dominates the uncertainty budget for
sea-level rise. Two independent multi-model exercises find deep
disagreement in Antarctica's own projected contribution. ISMIP6's
ensemble of process models disagrees even on its sign, spanning
outcomes from a small sea-level fall to 30~cm of rise by 2100 under
high emissions \citep{Seroussi2020}, while a sixteen-model synthesis
built on linear response theory finds a comparably wide spread in
basal-melt-driven mass loss \citep{Levermann2020}. A structured survey
of expert glaciological judgment, drawing on physical reasoning rather
than model runs, arrives independently at similarly wide, high-tailed
distributions \citep{Bamber2019}. At the emissions-scenario level,
Antarctica's own total contribution swings more dramatically with how
much warming is avoided than any other component does
\citep{DeConto2021}. A decomposition of total sea-level-rise variance
shows exactly when the balance tips. Emissions-scenario choice
dominates uncertainty through mid-century, after which geophysical
uncertainty, chiefly Antarctic ice dynamics, becomes the largest term
\citep{Darnell2025}. Under high emissions, the IPCC's own
low-likelihood, high-impact storyline for 2100 combines
earlier-than-projected ice-shelf disintegration, abrupt onset of these
Antarctic instabilities, and faster Greenland change to add more than a
meter of additional rise \citep{Fox-Kemper2021}. The near-term rate of
Antarctic mass loss even appears to predict its own rate three to five
decades ahead \citep{McCormack2026}, a relationship too newly reported
to narrow this range yet. This risk is structurally different from the
temperature-sensitive uncertainty in the other components. A single
central WAIS trajectory, rather than a mixture spanning its full range,
would understate exactly the high-end damages described at the outset.

Here we build a component-level framework that gives each contributor
to GMSLR its own parsimonious physical model, calibrated against an
independent observational record using Bayesian inference. The
contributors are ocean thermal expansion; Greenland Ice Sheet (GrIS)
surface mass balance and glacier discharge; West Antarctic Ice Sheet
(WAIS) glacier discharge; East Antarctic Ice Sheet (EAIS) surface mass
balance; the Antarctic Peninsula; glaciers outside the ice sheets; and
terrestrial water storage. We call it SPECTRUM, for Sea-level
Projections from Empirical Calibration of Thermodynamical,
Reduced-physics Models. Its only new component model links GrIS
discharge to ocean temperature through a 3-parameter time-lagged
response. Its treatment of WAIS is a structured deep-uncertainty
scenario mixture \citep{Bakker2017}, built to represent the skewed tail
described above rather than average over it. We withhold observed
GMSLR and its sensitivity to GMST, the relationship at the core of the
semi-empirical class described above, from every component calibration,
and use it only as a diagnostic on the composite model. The central
test is therefore out of sample. It asks whether the independently
calibrated parts reproduce a whole they were never fit to.

%% ====================================================================
\section*{Version 6 --- merged synthesis of Versions 3--5, WAIS transition corrected}

Global mean sea-level rise (GMSLR) already exposes roughly 1.6
million more people to annual flooding for every additional centimeter
\citep{Kulp2019}, with unmitigated damages projected at US\$14--18T
(trillion) per year once rise reaches 0.9--1.1~m \citep{Jevrejeva2018}.
Sea-level projections meant to inform that risk have been published for
four decades and span a wide range of methods and outcomes
\citep{Garner2026}. Rates of GMSLR have already doubled twice since the
mid-20th century \citep{Frederikse2020,Hamlington2024}. Coastal planning
nonetheless leans heavily on the IPCC Sixth Assessment Report (AR6),
which projects 0.4--0.7~m of 21st-century rise
\citep{Fox-Kemper2021,Hirschfeld2023}. A growing, cross-institutional
view holds that narrowing AR6's range requires tighter observational
constraint on projections, not more physical detail \citep{Felikson2025}.

Three model classes have taken up that challenge, each with a real
weakness. Process-based ice-sheet models resolve ice flow and
fracture directly, but they demand more constraints than the
observational record can currently supply. In hindcasting studies, they
diverge from the historical GMSLR record \citep{Aschwanden2021} and from
satellite observations within a decade of initialization
\citep{Fricker2025}. These disagreements trace to gaps such as an
assumed ice viscosity far less sensitive to stress than observations
indicate \citep{Millstein2022,GetraerMorlighem2025}. Forecasting
research spanning decades and many fields finds the same pattern.
Models with more free parameters than their data can constrain tend to
forecast worse than simpler alternatives \citep{Green2015,Makridakis2020}.
Semi-empirical models regress the rate of GMSLR onto global mean surface
temperature. They are simple and transparent, but the underlying
relationship is not obviously physical, causal, or stationary
\citep{Rahmstorf2007,Vermeer2009,Church2013}. Even so, such fits track
the observed record at least as well as contemporaneous process-based
projections \citep{Rahmstorf2012}. An even simpler benchmark needs no
temperature regression at all: a quadratic extrapolation of the
satellite-observed GMSL record alone agrees with IPCC SROCC and AR6
projections to within their own stated uncertainty \citep{Nerem2022}.
Between these extremes,
component-level models decompose GMSLR into its physical contributors
and calibrate each against existing data
\citep{Kopp2014,Mengel2016,Wong2017}, trading basin-scale resolution for
parsimony and a direct physical connection, while assuming the
relationships they calibrate are stationary. None of the published
implementations, however, withholds the aggregate observational record
from every component fit, so none can use it as an independent check on
the summed result. No one class is authoritative on its own. Ranked
instead by complexity, the three become a diagnostic. The simplest,
most transparent models motivate and check the more complex ones.

West Antarctica is the largest single source of uncertainty in
sea-level projections. That is not a consequence of how any model class
chooses to calibrate; it follows from the ice physics and the
observational record themselves. Four independent lines of evidence
converge on this conclusion, and West Antarctica's share of the total
uncertainty is growing.

First, the observational record shows the loss is already under way and
proceeding unevenly. Antarctica's grounding lines have retreated for
thirty consecutive years \citep{Rignot2026}, and the Thwaites Glacier
basin may already be undergoing a self-sustaining collapse
\citep{Joughin2014}. Radar surveys there find warm water undercutting
some sectors of the grounding line ten times faster than others
\citep{Milillo2019}, and recent observations remain consistent with
continued high rates of loss over the next fifty years
\citep{Goldberg2026}. Second, a structured survey of expert judgment,
built independently of any process-model ensemble, assigns West
Antarctica wider and more heavily tailed contribution distributions than
those ensembles do \citep{Bamber2019}. Third, that spread is not merely
disagreement between models talking past each other. A single ice-sheet
model's own uncertainty in calving, basal sliding, ice-shelf rheology,
and ocean forcing, propagated forward for centuries, already produces a
wide range of Antarctic outcomes \citep{Bulthuis2019}. The largest
multi-model ensemble, ISMIP6, disagrees even on the sign of Antarctica's
contribution \citep{Seroussi2020}, and a sixteen-model synthesis built
on linear response theory finds a comparably wide spread in
basal-melt-driven mass loss \citep{Levermann2020}. Fourth, that
disagreement is not fixed in time. A variance decomposition of
projected sea-level rise finds emissions-scenario choice dominant
through mid-century, after which geophysical uncertainty, chiefly
Antarctic ice dynamics, becomes the largest term \citep{Darnell2025},
and most of the emissions-pathway spread in total projections traces to
the Antarctic contribution alone \citep{DeConto2021}. Under high
emissions, the IPCC's own low-likelihood, high-impact storyline for
2100 combines earlier-than-projected ice-shelf disintegration, abrupt
onset of marine ice sheet instability, and faster Greenland change to
add more than a meter of additional rise \citep{Fox-Kemper2021}. One
recent finding cuts the other way: the near-term rate of Antarctic mass
loss appears to carry predictive information about the rate over the
following three to five decades \citep{McCormack2026}, though this
relationship is too newly reported to narrow the range yet.

The physical mechanism behind West Antarctic loss argues against
reading that predictive signal as reassurance. Marine ice sheet
instability does not merely widen projection uncertainty; it skews the
distribution toward high outcomes \citep{Robel2019}. Coupled with
hydrofracturing and ice-cliff collapse, that instability could add more
than a meter of Antarctic sea-level rise by 2100 under unabated
emissions \citep{DeContoP2016}. How much of that process will actually
occur is itself disputed: recalibrating the same physics against
observations from the last interglacial, the Pliocene, and the past 25
years finds no need to invoke ice-cliff collapse at all, bounding the
95th-percentile Antarctic contribution below 43~cm this century without
it \citep{Edwards2019}. Which physical processes to assume, not only
which parameters to fit, therefore changes the shape of West
Antarctica's projected distribution \citep{Kopp2017}. Independent
regional ocean modeling adds a further constraint: shelf melt
accelerates at similar rates across emissions scenarios this century,
so much of this risk is already locked in regardless of the mitigation
path taken \citep{Naughten2023}. This uncertainty is structurally
different from the temperature-sensitive uncertainty in the other
components, which is why some published projections already treat West
Antarctica as a deep-uncertainty mixture of scenarios rather than a
single central trajectory \citep{Bakker2017}. Anything less would
understate exactly the high-end damages introduced above.

No published component-level framework performs the independent check
described above, and the gap matters most exactly where the stakes are
largest: West Antarctica. Here, we build a framework in which each
contributor to GMSLR gets its own parsimonious physical model,
calibrated against an independent observational record using Bayesian
inference. The contributors are ocean thermal expansion; Greenland Ice
Sheet (GrIS) surface mass balance and glacier discharge; West Antarctic
Ice Sheet (WAIS) glacier discharge; East Antarctic Ice Sheet (EAIS)
surface mass balance; the Antarctic Peninsula; glaciers outside the ice
sheets; and terrestrial water storage. We call it SPECTRUM, for
Sea-level Projections from Empirical Calibration of Thermodynamical,
Reduced-physics Models. Its only new component model links GrIS
discharge to ocean temperature through a 3-parameter time-lagged
response. Its treatment of WAIS is a structured deep-uncertainty
scenario mixture, built to represent the skewed tail described above
rather than average over it. We withhold observed GMSLR and its
sensitivity to GMST, the relationship at the core of the semi-empirical
class described above, from every component calibration, and use it
only as a diagnostic on the composite model. The central test is
therefore out of sample: it asks whether the independently calibrated
parts reproduce a whole they were never fit to.

\bibliographystyle{agufull08}
\bibliography{references}

\end{document}
